\documentclass[11pt, letterpaper]{article}

\usepackage[top=.8in, bottom=.8in, left=0.5in, right=0.5in]{geometry}
\usepackage{amsmath, amssymb, amsthm, mathtools}
\usepackage{longtable}
\usepackage{booktabs}
\usepackage{array}
\usepackage{tabularx}
\usepackage{xcolor}
\usepackage{fancyhdr}
\usepackage{titlesec}
\usepackage{setspace}
\usepackage[T1]{fontenc}
\usepackage{enumitem}
\usepackage{tocloft}
\usepackage{tcolorbox}
\tcbuselibrary{breakable, skins}
\usepackage{listings}
\usepackage{courier}

\definecolor{RSblue}{RGB}{25, 60, 120}
\definecolor{RSgray}{RGB}{90, 90, 90}
\definecolor{RSgreen}{RGB}{20, 110, 60}
\definecolor{RSred}{RGB}{160, 30, 30}
\definecolor{RSamber}{RGB}{180, 110, 0}
\definecolor{RSpurple}{RGB}{90, 30, 120}

\newtcolorbox{rsbox}{
  enhanced, breakable,
  colback=RSblue!4, colframe=RSblue!60,
  title={\textbf{RS Ontological Statement}},
  fonttitle=\small\bfseries\color{white},
  attach boxed title to top left={yshift=-2mm, xshift=4mm},
  boxed title style={colback=RSblue},
  left=4pt, right=4pt, top=6pt, bottom=4pt
}

\newtcolorbox{verifybox}{
  enhanced, breakable,
  colback=RSgreen!4, colframe=RSgreen!50,
  title={\textbf{External Verification (Continuous Lens)}},
  fonttitle=\small\bfseries\color{white},
  attach boxed title to top left={yshift=-2mm, xshift=4mm},
  boxed title style={colback=RSgreen!70!black},
  left=4pt, right=4pt, top=6pt, bottom=4pt
}

\newtcolorbox{observebox}{
  enhanced, breakable,
  colback=RSamber!5, colframe=RSamber!60,
  title={\textbf{Structural Observation}},
  fonttitle=\small\bfseries\color{white},
  attach boxed title to top left={yshift=-2mm, xshift=4mm},
  boxed title style={colback=RSamber!80!black},
  left=4pt, right=4pt, top=6pt, bottom=4pt
}


\lstset{
  basicstyle=\footnotesize\ttfamily,
  breaklines=true, frame=single,
  language=Python, showstringspaces=false,
  numbers=left, numberstyle=\tiny,
  commentstyle=\color{RSgray},
  keywordstyle=\color{RSblue!70!black}}

\pagestyle{fancy}
\fancyhf{}
\fancyhead[L]{\small\textcolor{RSgray}{RigbySpace Discrete Dynamics}}
\fancyhead[R]{\small\textcolor{RSgray}{Mediant Geodesics - Page \thepage}}
\fancyfoot[L]{\small\textcolor{RSgray}{D. Veneziano}}
\fancyfoot[R]{\small\textcolor{RSgray}{Discrete Unification Framework}}
\renewcommand{\headrulewidth}{0.4pt}

\titleformat{\section}
  {\large\bfseries\color{RSblue}}{\thesection}{1em}{}[\titlerule]
\titleformat{\subsection}
  {\normalsize\bfseries\color{RSblue}}{\thesubsection}{1em}{}
\titleformat{\subsubsection}
  {\normalsize\bfseries\color{RSgray}}{\thesubsubsection}{1em}{}

\onehalfspacing
\setlength{\parskip}{4pt}
\setlength{\parindent}{0pt}

\newcolumntype{L}[1]{>{\raggedright\arraybackslash}p{#1}}

\newcommand{\Z}{\mathbb{Z}}
\newcommand{\Sinter}{S_{\mathrm{int}}}
\newcommand{\SL}{S_{\mathrm{L}}}
\newcommand{\SLe}{S_{\mathrm{L}}^{\mathrm{ext}}}
\newcommand{\SP}{S_{\mathrm{P}}}
\newcommand{\SPs}{S_{\mathrm{P}}^{*}}
\newcommand{\SPe}{S_{\mathrm{P}}^{\mathrm{ext}}}
\newcommand{\IP}{I_{\mathrm{P}}}
\newcommand{\Vac}{\mathcal{V}}
\newcommand{\SLpos}{S_{\mathrm{L}}^{+}}
\newcommand{\GK}{\Gamma_{K}}
\newcommand{\Ethree}{E_{3+1}}
\newcommand{\SRS}{S_{\mathrm{RS}}}
\newcommand{\TRS}{\mathcal{T}_{\mathrm{RS}}}
\newcommand{\Ens}{\mathcal{E}}

\newcommand{\lop}{\lambda}
\newcommand{\aop}{\eta}

\newcommand{\bplus}{\boxplus}
\newcommand{\med}{\oplus}
\newcommand{\ecplus}{\oplus_{\mathrm{ec}}}
\renewcommand{\succ}{\operatorname{succ}}

\newcommand{\rank}{\rho}
\newcommand{\hfn}{h}
\newcommand{\Dcross}{\Delta_{\mathrm{cross}}}
\newcommand{\sgn}{\operatorname{sgn}}
\newcommand{\rem}{\operatorname{rem}}

\newcommand{\koppa}{\mathbin{\text{\rm\k{o}}}}

\newcommand{\Tcyc}{T_{11}}
\newcommand{\PhiSet}{\Phi}
\newcommand{\Tbary}{T_{\mathrm{bary}}}
\newcommand{\Psistate}{\Psi}
\newcommand{\Tcharge}{\Gamma}

\newcommand{\Dg}{D_{\mathrm{gauge}}}
\newcommand{\Sint}{S_{\mathrm{int}}}
\newcommand{\Bridge}{\mathcal{B}}
\newcommand{\Lalpha}{\alpha_{\mathrm{lat}}}
\newcommand{\LQ}{\mathcal{L}_{\mathbb{Z}}}
\newcommand{\Sigimb}{\Sigma_{\mathrm{imb}}}
\newcommand{\VS}{\mathcal{V}_S}
\newcommand{\PG}{\mathrm{PG}}

% --- New macros: RS Sommerfeld and Pi-Analog ---
\newcommand{\Gmg}{G_{\mathrm{mg}}}
\newcommand{\Ksat}{K_{\mathrm{sat}}}
\newcommand{\dslip}{\delta_{\mathrm{slip}}}
\newcommand{\Tvac}{T_{\mathrm{vac}}}
\newcommand{\tauX}{\tau_{10}}
\newcommand{\Pithree}{\Pi_3}
\newcommand{\nG}{n_G}
\newcommand{\aSC}{\alpha^{-1}_{\mathrm{int}}}
\newcommand{\piRS}{\pi_{\mathrm{RS}}}
\newcommand{\pa}{\mathfrak{p}_\alpha}
\newcommand{\qa}{\mathfrak{q}_\alpha}
\newcommand{\Da}{\mathcal{D}_\alpha}
\newcommand{\qosc}{\mathfrak{q}_{\mathrm{osc}}}

\theoremstyle{definition}
\newtheorem{definition}{Definition}[section]
\newtheorem{axiom}{Axiom}[section]
\newtheorem{remark}{Remark}[section]
\newtheorem{postulate}{Postulate}[section]
\newtheorem{law}{Law}[section]
\newtheorem{observation}{Observation}[section]

\theoremstyle{plain}
\newtheorem{theorem}{Theorem}[section]
\newtheorem{lemma}{Lemma}[section]
\newtheorem{proposition}{Proposition}[section]
\newtheorem{corollary}{Corollary}[section]



\begin{document}

\begin{titlepage}
  \centering
  \vspace*{2cm}
  {\Huge\bfseries\color{RSblue} RigbySpace Discrete Dynamics\par}
  \vspace{0.5cm}
  {\LARGE Mediant Geodesics, Saturation, and Boundary Transport\par}
  \vspace{0.5cm}
  {\large RigbySpace Discrete Integer Profile\par}
  \vspace{2cm}
  \rule{0.6\textwidth}{0.4pt}\par
  \vspace{1cm}
  {\large D.\ Veneziano\par}
  \vspace{0.8cm}
  {\normalsize March, 2026\par}
  \vspace{1.5cm}
  \rule{0.6\textwidth}{0.4pt}\par
  \vspace{0.5cm}
  {\small Propagation, adjacency, saturation, and coupled endpoint transport\par}
  \vfill
  {\small\color{RSgray}
  This document adopts the current master reference as authoritative.  Every
  definition and derivation below is written inside the RigbySpace ontology and
  uses the master preamble, colour system, and box structures.\par}
\end{titlepage}

\tableofcontents
\newpage

\section{Scope and primitive starting set}

This document replaces the earlier mediant-geodesic draft with a version that is
conformant with the current master reference.  The present document keeps the
subject matter of mediant trees, geodesic paths, saturation, and endpoint
transport, while removing statements that depended on continuum assumptions,
single-state interpretations of $\psi$, or non-master notions of height and
adjacency.

The derivation begins from the RS primitives before any numerical family is
introduced.  The operative primitives for the present topic are the mediant
sum $\med$, the transforms $\lop$ and $\aop$, the coupled reciprocal $\psi$,
the ledger operator $\koppa$, the cycle boundary $\Omega$, and the master
ordering/width structure built from $\Dcross$.  The global RS constants
$\aSC$ and $\piRS$ belong to the larger ontology and remain available as
master primitives, but they are not required for the mediant-geodesic chain
itself.

\begin{rsbox}
The mediant-geodesic sector is a structural sector of the master reference.  It
does not create a second ontology.  Its admissible content is therefore limited
to integer pairs, cross-ordering, width, mediant insertion, saturation,
barycentric oscillation, coupled reciprocal transport, ZERO-context boundary
states, and ledger selection at $\Omega$.
\end{rsbox}

\begin{definition}[Explicit Rational Pair (ERP)]
An ERP is an ordered pair $(n,d)\in\Z\times(\Z\setminus\{0\})$.  The numerator
encodes event magnitude and the denominator encodes geometric capacity.  ERPs
are never reduced to lowest terms.
\end{definition}

\begin{definition}[Observational equivalence]
For $A=(n_1,d_1)$ and $B=(n_2,d_2)$ in $\SL$, one writes $A\sim B$ when
$n_1 d_2=n_2 d_1$.  This is the only admissible equivalence on ERPs.  GCD
reduction is not a rule of the ontology.
\end{definition}

\begin{definition}[Cross-ordering and width]
For $A=(p_a,q_a)$ and $B=(p_b,q_b)$ in $\SL$, define
\[
\Dcross(A,B):=p_b q_a-p_a q_b.
\]
Write $A\prec B$ when $\Dcross(A,B)>0$.  If $L\prec U$, define the width
\[
W(L,U):=\Dcross(L,U)\in\Z_{>0}.
\]
The bracket is saturated when $W(L,U)=1$.
\end{definition}

\begin{definition}[Mediant sum]
For $A=(a,b)$ and $B=(c,d)$ in $\SL$, define
\[
A\med B:=(a+c,b+d).
\]
The mediant sum is an ERP-valued operation internal to the integer lattice.
\end{definition}

\begin{definition}[The transforms $\lop$ and $\aop$]
The master transforms are
\[
\lop(n,d):=(n+d,d),
\qquad
\aop(n,d):=(n+d,n).
\]
The transform $\lop$ preserves the denominator while increasing the numerator
by one denominator-load.  The transform $\aop$ turns the old numerator into the
new denominator and is the master Fibonacci generator.
\end{definition}

\begin{definition}[Coupled reciprocal $\psi$]
For a coupled pair of rational states,
\[
\psi\!\left(\frac{a}{b},\,\frac{c}{d}\right)
=
\left(\frac{d}{a},\,\frac{b}{c}\right).
\]
The operation is coupled.  It does not act on a single ERP in isolation.  Its
second iterate exchanges the two states and its fourth iterate is the identity.
\end{definition}

\begin{definition}[Boundary seeds and ZERO-context roles]
The mediant-tree root bracket uses the formal seeds $(0,1)$ and $(1,0)$ from
the master reference.  These are not asserted here as ordinary active matter
states.  The label $(0,1)$ is the empty-position label appropriate to mediant
propagation context.  The label $(1,0)$ is a constraint-vacuum boundary seed
with absent metric basis.  Their use in the tree is therefore contextual and
fully compatible with the master ZERO logic.
\end{definition}

\begin{theorem}[Width preservation under mediant insertion]
\begin{rsbox}
Let $L=(a,b)$ and $U=(c,d)$ with $L\prec U$, and let $M:=L\med U=(a+c,b+d)$.
Then
\[
\Dcross(L,M)=\Dcross(L,U),
\qquad
\Dcross(M,U)=\Dcross(L,U).
\]
Consequently,
\[
W(L,M)=W(M,U)=W(L,U).
\]
\end{rsbox}
\begin{proof}
\begin{rsbox}
A direct computation gives
\[
\Dcross(L,M)=(a+c)b-a(b+d)=cb-ad=\Dcross(L,U),
\]
and
\[
\Dcross(M,U)=c(b+d)-(a+c)d=cb-ad=\Dcross(L,U).
\]
Hence both child brackets inherit the same width as the parent bracket.
\end{rsbox}
\end{proof}

\end{theorem}

\begin{theorem}[The $\lop$-ray from the unit ERP is saturated]
\begin{rsbox}
For every integer $k\geq 0$,
\[
\lop^{k}(1,1)=(k+1,1).
\]
Moreover,
\[
\Dcross\!\bigl(\lop^{k}(1,1),\lop^{k+1}(1,1)\bigr)=1,
\]
so consecutive states on the $\lop$-ray form a saturated chain.
\end{rsbox}
\begin{proof}
\begin{rsbox}
The identity $\lop(n,d)=(n+d,d)$ implies by induction that
$\lop^{k}(1,1)=(k+1,1)$.  Therefore
\[
\Dcross\!\bigl((k+1,1),(k+2,1)\bigr)=(k+2)\cdot 1-(k+1)\cdot 1=1.
\]
Thus every consecutive pair is adjacent in the saturated sense.
\end{rsbox}
\end{proof}

\end{theorem}

\begin{theorem}[The $\aop$-ray from the unit ERP generates the Fibonacci edge chain]
\begin{rsbox}
Let $F_1:=1$, $F_2:=1$, and $F_{m+2}:=F_{m+1}+F_m$ for $m\geq 1$.  Then for
every integer $k\geq 0$,
\[
\aop^{k}(1,1)=\bigl(F_{k+2},F_{k+1}\bigr).
\]
In addition,
\[
\Dcross\!\bigl(\aop^{k}(1,1),\aop^{k+1}(1,1)\bigr)
=
\begin{cases}
1 & \text{if }k\text{ is even},\\
-1 & \text{if }k\text{ is odd}.
\end{cases}
\]
\end{rsbox}
\begin{proof}
\begin{rsbox}
The statement is true for $k=0$ because $\aop^0(1,1)=(1,1)=(F_2,F_1)$.  Assume
$\aop^k(1,1)=(F_{k+2},F_{k+1})$.  Then
\[
\aop^{k+1}(1,1)=\aop(F_{k+2},F_{k+1})
=(F_{k+3},F_{k+2}),
\]
which proves the Fibonacci formula by induction.  For the cross-determinant,
\[
\Dcross\!\bigl((F_{k+2},F_{k+1}),(F_{k+3},F_{k+2})\bigr)
=
F_{k+3}F_{k+1}-F_{k+2}^2.
\]
Cassini's identity yields
\[
F_{k+3}F_{k+1}-F_{k+2}^2=(-1)^k.
\]
Writing this by parity gives the displayed formula.
\end{rsbox}
\end{proof}

\end{theorem}

\begin{observebox}
The two previous theorems satisfy the operator-first requirement.  The
numerical families appear only after the transforms $\lop$ and $\aop$ have
already fixed the state update law.  The $\lop$-ray supplies a monotone
capacity-preserving saturated edge chain.  The $\aop$-ray supplies the canonical
recurrence edge chain associated with the master Fibonacci generator.
\end{observebox}

\section{Ordered brackets and the mediant tree}

\begin{definition}[Ordered boundary bracket]
An ordered boundary bracket is a pair $[L,U]$ of ERPs with $L\prec U$.  Its
interior mediator is the ERP $L\med U$.  A bracket is saturated when
$W(L,U)=1$.
\end{definition}

\begin{definition}[Root bracket]
The root bracket of the mediant tree is
\[
B_0:=[(0,1),(1,0)].
\]
Its first interior state is the mediant
\[
(0,1)\med(1,0)=(1,1).
\]
\end{definition}

\begin{proposition}[The root bracket is saturated]
\begin{rsbox}
The root bracket $B_0=[(0,1),(1,0)]$ has width one.
\end{rsbox}
\begin{proof}
\begin{rsbox}
By direct calculation,
\[
W((0,1),(1,0))=\Dcross((0,1),(1,0))=1\cdot 1-0\cdot 0=1.
\]
Hence the root bracket is saturated.
\end{rsbox}
\end{proof}

\end{proposition}

\begin{definition}[Left and right mediant refinement]
For an ordered bracket $[L,U]$, define its two child brackets by
\[
\mathcal{R}_{\mathrm{L}}[L,U]:=[L,L\med U],
\qquad
\mathcal{R}_{\mathrm{R}}[L,U]:=[L\med U,U].
\]
These are the elementary refinement operations of the mediant tree.
\end{definition}

\begin{proposition}[Refinement preserves ordering]
\begin{rsbox}
Let $[L,U]$ be an ordered bracket and let $M=L\med U$.  Then
\[
L\prec M\prec U.
\]
\end{rsbox}
\begin{proof}
\begin{rsbox}
Because $L\prec U$, the previous width-preservation theorem gives
\[
\Dcross(L,M)=\Dcross(L,U)>0
\]
and
\[
\Dcross(M,U)=\Dcross(L,U)>0.
\]
Hence $L\prec M$ and $M\prec U$.
\end{rsbox}
\end{proof}

\end{proposition}

\begin{definition}[Mediant tree $\mathcal{T}$]
The mediant tree is the rooted binary tree obtained by recursive application of
the two refinement operations to the root bracket $B_0$.  Every refinement
creates one new interior ERP and two new ordered boundary brackets.  The
vertices of the tree are the ERPs produced in this recursion, and two vertices
share an edge precisely when they appear as consecutive vertices on one parent
to child branch of the recursion.
\end{definition}

\begin{theorem}[Unique bracket ancestry]
\begin{rsbox}
Every vertex of $\mathcal{T}$ is created by a unique finite word in the two
refinement operations $\mathcal{R}_{\mathrm{L}}$ and
$\mathcal{R}_{\mathrm{R}}$ applied to the root bracket.
\end{rsbox}
\begin{proof}
\begin{rsbox}
At each stage the recursion replaces one bracket by two children whose interiors
lie in disjoint ordered positions relative to the parent mediator.  A new
interior ERP is therefore created in exactly one child bracket at that stage.
Proceeding inductively from the root, the finite sequence of left and right
choices that creates a given vertex is unique.
\end{rsbox}
\end{proof}

\end{theorem}

\begin{theorem}[Unique path property]
\begin{rsbox}
Between any two vertices $A$ and $B$ of $\mathcal{T}$ there exists exactly one
simple path in the graph of the tree.
\end{rsbox}
\begin{proof}
\begin{rsbox}
The recursive construction produces a rooted binary tree with no cycles.  In
every tree, the path from $A$ to $B$ is obtained by ascending from $A$ to the
lowest common ancestor of $A$ and $B$ and then descending to $B$.  Because the
ancestor chain of each vertex is unique, this path is unique.
\end{rsbox}
\end{proof}

\end{theorem}

\begin{observebox}
The uniqueness statement is a structural statement about the mediant tree as a
tree.  It does not depend on GCD reduction, Euclidean simplification, or any
continuum identification of rationals.  The object carrying the path is the ERP
vertex together with its place in the recursive refinement graph.
\end{observebox}

\section{Mediant geodesics and structural symbols}

\begin{definition}[Mediant geodesic $\gamma(A,B)$]
For vertices $A$ and $B$ of $\mathcal{T}$, the mediant geodesic $\gamma(A,B)$
is the unique simple path in $\mathcal{T}$ from $A$ to $B$.
\end{definition}

\begin{proposition}[Existence and uniqueness of mediant geodesics]
\begin{rsbox}
For any vertices $A$ and $B$ of $\mathcal{T}$, the mediant geodesic
$\gamma(A,B)$ exists and is unique.
\end{rsbox}
\begin{proof}
\begin{rsbox}
This is the unique path property of the tree stated in geodesic language.
\end{rsbox}
\end{proof}

\end{proposition}

\begin{definition}[Geodesic length]
If
\[
\gamma(A,B)=(S_0,S_1,\dots,S_m)
\]
with $S_0=A$ and $S_m=B$, then the geodesic length is
\[
\ell(\gamma(A,B)):=m.
\]
\end{definition}

\begin{lemma}[Concatenation through an intermediate vertex]
\begin{rsbox}
Let $A$, $B$, and $C$ be vertices of $\mathcal{T}$, and suppose $B$ lies on the
geodesic $\gamma(A,C)$.  Then
\[
\gamma(A,C)=\gamma(A,B)\cup\gamma(B,C)
\]
with the two subpaths meeting only at the vertex $B$.
\end{rsbox}
\begin{proof}
\begin{rsbox}
Because $\gamma(A,C)$ is the unique simple path from $A$ to $C$, any vertex
lying on it divides the path into an initial segment and a terminal segment.
Those two segments are respectively the unique simple paths from $A$ to $B$ and
from $B$ to $C$.
\end{rsbox}
\end{proof}

\end{lemma}

\begin{corollary}[Length additivity]
\begin{rsbox}
Under the hypotheses of the preceding lemma,
\[
\ell(\gamma(A,C))=\ell(\gamma(A,B))+\ell(\gamma(B,C)).
\]
\end{rsbox}
\begin{proof}
\begin{rsbox}
The path $\gamma(A,C)$ is the concatenation of two edge-disjoint geodesic
segments sharing only the vertex $B$, so the edge counts add.
\end{rsbox}
\end{proof}

\end{corollary}

\begin{definition}[Structural symbol $\{A,B\}$]
The structural symbol $\{A,B\}$ is the formal sum of the oriented mediant
geodesic edges from $A$ to $B$.  The defining relations are
\[
\{A,B\}+\{B,C\}=\{A,C\},
\qquad
\{A,A\}=0,
\qquad
\{A,B\}+\{B,A\}=0.
\]
\end{definition}

\begin{theorem}[Structural symbol additivity]
\begin{rsbox}
If $B$ lies on $\gamma(A,C)$, then
\[
\{A,C\}=\{A,B\}+\{B,C\}.
\]
\end{rsbox}
\begin{proof}
\begin{rsbox}
The formal oriented edge sum for $\gamma(A,C)$ is exactly the disjoint
concatenation of the oriented edge sums for $\gamma(A,B)$ and $\gamma(B,C)$.
\end{rsbox}
\end{proof}

\end{theorem}

\begin{corollary}[Closed triangular cancellation]
\begin{rsbox}
For any vertices $A$, $B$, and $C$,
\[
\{A,B\}+\{B,C\}+\{C,A\}=0.
\]
\end{rsbox}
\begin{proof}
\begin{rsbox}
By additivity,
\[
\{A,B\}+\{B,C\}=\{A,C\}.
\]
Adding $\{C,A\}$ to both sides and using $\{A,C\}+\{C,A\}=0$ gives the result.
\end{rsbox}
\end{proof}

\end{corollary}

\begin{definition}[Atomic edge symbol]
If $L\prec U$ and $W(L,U)=1$, then the geodesic $\gamma(L,U)$ has one edge and
the structural symbol $\{L,U\}$ is called an atomic edge symbol.
\end{definition}

\begin{proposition}[Mediant insertion splits an atomic edge into two atomic edges]
\begin{rsbox}
Let $L\prec U$ with $W(L,U)=1$, and let $M:=L\med U$.  Then
\[
\{L,U\}=\{L,M\}+\{M,U\},
\]
and both $\{L,M\}$ and $\{M,U\}$ are atomic edge symbols.
\end{rsbox}
\begin{proof}
\begin{rsbox}
Width preservation gives
\[
W(L,M)=W(M,U)=W(L,U)=1.
\]
Hence both child brackets are saturated, and the structural symbol additivity
relation yields the decomposition.
\end{rsbox}
\end{proof}

\end{proposition}

\section{Saturation, barycentric oscillation, and selection at $\Omega$}

\begin{definition}[Barycentric oscillation]
A barycentric oscillation is the stable two-state regime associated with a
saturated bracket $[L,U]$ of width one.  In this regime the ontology does not
replace the pair by a continuum point.  The physically admissible attractor is
the ordered alternation between the adjacent ERPs $L$ and $U$.
\end{definition}

\begin{proposition}[Saturation is not collapse]
\begin{rsbox}
Let $[L,U]$ be saturated and set $M:=L\med U$.  Then the child brackets
$[L,M]$ and $[M,U]$ are also saturated.  Therefore the saturated regime is not
a collapse to a single state but a persistent width-one bracket structure.
\end{rsbox}
\begin{proof}
\begin{rsbox}
This is immediate from width preservation.  Saturation is stable under mediant
insertion.  The ontology therefore records a permanent adjacent-state pair
rather than a real-valued point limit.
\end{rsbox}
\end{proof}

\end{proposition}

\begin{definition}[Barycentric truncation representative]
For a saturated bracket $[L,U]$, define the selected representative
\[
S_{\mathrm{rep}}=\mathcal{T}_B([L,U],\koppa_{\mathrm{ledger}})
\]
by the master rule
\[
S_{\mathrm{rep}}:=
\begin{cases}
L & \text{if }\rank(\koppa_{\mathrm{ledger}})\text{ is even},\\
U & \text{if }\rank(\koppa_{\mathrm{ledger}})\text{ is odd}.
\end{cases}
\]
This representative is chosen at the cycle boundary $\Omega$.
\end{definition}

\begin{theorem}[Parity-stable selection law]
\begin{rsbox}
Let $[L,U]$ be saturated.  If two consecutive evaluations at $\Omega$ are made
with ledger values whose structural ranks have the same parity, then the
selected representative is the same at both evaluations.  If the parity flips,
the representative swaps from one endpoint of the bracket to the other.
\end{rsbox}
\begin{proof}
\begin{rsbox}
The selection rule depends only on the parity class of
$\rank(\koppa_{\mathrm{ledger}})$.  Equal parity yields the same branch of the
definition, and opposite parity yields the other branch.
\end{rsbox}
\end{proof}

\end{theorem}

\begin{observebox}
This is the correct RS replacement for continuum convergence language.  A
saturated bracket carries two adjacent states, and the ledger at $\Omega$
chooses which state is exposed as the representative state for the next cycle.
No real-number limit is introduced anywhere in the chain.
\end{observebox}

\section{Coupled transport, orientation reversal, and ZERO-boundary behaviour}

\begin{definition}[Endpoint-paired $\psi$ transport]
Let
\[
A=\frac{a}{b},
\qquad
B=\frac{c}{d}
\]
be a coupled pair of endpoint states associated with an oriented geodesic datum.
Their $\psi$-transport is
\[
\Psistate(A,B):=
\psi\!\left(\frac{a}{b},\frac{c}{d}\right)
=
\left(\frac{d}{a},\frac{b}{c}\right).
\]
\end{definition}

\begin{theorem}[$\psi^2$ reverses geodesic orientation]
\begin{rsbox}
For any coupled endpoint pair $(A,B)$,
\[
\psi^2(A,B)=(B,A).
\]
Consequently,
\[
\{(\psi^2(A,B))_1,(\psi^2(A,B))_2\}=\{B,A\}=-\{A,B\}.
\]
\end{rsbox}
\begin{proof}
\begin{rsbox}
The master definition of $\psi$ gives
\[
\psi^2\!\left(\frac{a}{b},\frac{c}{d}\right)
=
\left(\frac{c}{d},\frac{a}{b}\right).
\]
Thus the second iterate exchanges the endpoint order.  The structural symbol
relation $\{A,B\}+\{B,A\}=0$ then yields
\[
\{B,A\}=-\{A,B\}.
\]
\end{rsbox}
\end{proof}

\end{theorem}

\begin{definition}[Constraint-vacuum endpoint]
An endpoint label of the form $n/0$ with $n\neq 0$ is a constraint vacuum.
Its structural history is preserved while the metric basis is absent.  Such a
state cannot continue as an active endpoint until $\psi$ resolves the
suspension through pairing with an active partner.
\end{definition}

\begin{proposition}[ZERO-boundary resolution of a suspended endpoint]
\begin{rsbox}
Let one endpoint be suspended and the other active:
\[
\left(\frac{n}{0},\frac{c}{d}\right).
\]
Then one $\psi$ step yields
\[
\psi\!\left(\frac{n}{0},\frac{c}{d}\right)
=
\left(\frac{d}{n},\frac{0}{c}\right).
\]
The first output is an active resolved state and the second output is an
empty-position trace carrying contextual history in its denominator.
\end{rsbox}
\begin{proof}
\begin{rsbox}
This is the master ZERO-resolution formula applied directly to the endpoint
pair.  The history-bearing numerator $n$ becomes the denominator of the active
resolved state $d/n$, while the partner magnitude $c$ becomes the denominator of
the zero-magnitude trace $0/c$.
\end{rsbox}
\end{proof}

\end{proposition}

\begin{observebox}
The lower root seed $(0,1)$ belongs to the same ZERO-context family as the
zero-magnitude labels.  It is a propagation-context label for emptiness, not an
ordinary matter state.  The mediant tree may therefore use the root bracket
$[(0,1),(1,0)]$ without collapsing ZERO logic into numerical zero.
\end{observebox}

\section{External verification}

\begin{verifybox}
The formal content above is ontological and integer-native.  The following
verification statements do not add ontology.  They only check the internal
claims by computation.

A verification program can construct the mediant tree by recursive refinement
from the root bracket $[(0,1),(1,0)]$, locate a target ERP by repeated
cross-order comparison, and build geodesics by ancestor comparison.  It can
check width preservation, uniqueness of geodesic paths, structural-symbol
concatenation, the $\lop$-ray formula, the $\aop$-ray Fibonacci formula,
barycentric truncation parity selection, and ZERO-boundary $\psi$ resolution.

The program included in the appendix does exactly this.  It avoids GCD
reduction in the core tree logic, uses observational equivalence for endpoint
matching, and implements $\koppa$ by iterative subtraction.
\end{verifybox}

\appendix

\section{Python implementation and verification appendix}

The appendix implements every mathematical construct introduced in this document:
ERP data, observational equivalence, $\Dcross$, ordering, width, mediant
insertion, the transforms $\lop$ and $\aop$, coupled $\psi$ transport,
constructive remainder $\koppa$, structural rank, root-bracket refinement,
target location in the tree, geodesic extraction, structural symbols as edge
lists, barycentric truncation, and ZERO-boundary resolution.  The test suite
checks the document's core statements.

\begin{lstlisting}
from dataclasses import dataclass

G_MG = 2

@dataclass(frozen=True)
class ERP:
    n: int
    d: int

def obs_equiv(a: ERP, b: ERP) -> bool:
    return a.n * b.d == b.n * a.d

def d_cross(a: ERP, b: ERP) -> int:
    return b.n * a.d - a.n * b.d

def prec(a: ERP, b: ERP) -> bool:
    return d_cross(a, b) > 0

def width(l: ERP, u: ERP) -> int:
    if not prec(l, u):
        raise ValueError("Bracket is not ordered")
    return d_cross(l, u)

def mediant(a: ERP, b: ERP) -> ERP:
    return ERP(a.n + b.n, a.d + b.d)

def lambda_transform(x: ERP) -> ERP:
    return ERP(x.n + x.d, x.d)

def eta_transform(x: ERP) -> ERP:
    return ERP(x.n + x.d, x.n)

def psi_pair(a: ERP, b: ERP):
    return ERP(b.d, a.n), ERP(a.d, b.n)

def zero_resolution(n: int, c: int, d: int):
    return ERP(d, n), ERP(0, c)

def koppa(x: int, y: int) -> int:
    if y <= 0:
        raise ValueError("Second argument must be positive")
    r = x
    while r >= y:
        r -= y
    while r < 0:
        r += y
    return r

def rank_value(m: int) -> int:
    # Structural rank: total prime factors counted with multiplicity.
    # Primality of p detected by koppa: p is prime iff koppa(p, k) != 0
    # for all k with 1 < k < p. Division used here is in the verification
    # (continuous-lens) context only.
    if m <= 1:
        return 0
    total = 0
    p = 2
    while p <= m:
        # Check primality of p using koppa
        is_prime = all(koppa(p, k) != 0 for k in range(2, p))
        if not is_prime:
            p += 1
            continue
        while koppa(m, p) == 0:
            total += 1
            # Divide out one factor of p (verification context)
            q, r = divmod(m, p)
            assert r == 0
            m = q
        p += 1
    return total

def lop_ray(k: int):
    x = ERP(1, 1)
    for _ in range(k):
        x = lambda_transform(x)
    return x

def eta_ray(k: int):
    x = ERP(1, 1)
    for _ in range(k):
        x = eta_transform(x)
    return x

def locate_target(target: ERP, max_steps: int = 256):
    # Boundary seeds (0,1) and (1,0) are the tree roots, not interior vertices.
    # They are not reachable by the mediant search loop and are returned directly
    # with an empty ancestry path.
    if obs_equiv(target, ERP(0, 1)) or obs_equiv(target, ERP(1, 0)):
        return [target]
    left = ERP(0, 1)
    right = ERP(1, 0)
    path = [mediant(left, right)]
    steps = 0
    while steps < max_steps:
        mid = mediant(left, right)
        if obs_equiv(mid, target):
            return path
        if prec(left, target) and prec(target, mid):
            right = mid
        elif prec(mid, target) and prec(target, right):
            left = mid
        else:
            raise ValueError("Target outside current bracket")
        mid = mediant(left, right)
        path.append(mid)
        steps += 1
    raise RuntimeError("Target not found within step limit")

def geodesic(a: ERP, b: ERP):
    path_a = locate_target(a)
    path_b = locate_target(b)
    i = 0
    while i < min(len(path_a), len(path_b)) and obs_equiv(path_a[i], path_b[i]):
        i += 1
    lca = i - 1
    left_part = list(reversed(path_a[lca+1:]))
    right_part = path_b[lca:]
    return left_part + right_part

def structural_symbol(a: ERP, b: ERP):
    g = geodesic(a, b)
    return list(zip(g[:-1], g[1:]))

def barycentric_truncation(l: ERP, u: ERP, ledger_value: int) -> ERP:
    if rank_value(ledger_value) % 2 == 0:
        return l
    return u

def check_width_preservation():
    l = ERP(1, 3)
    u = ERP(2, 5)
    if not prec(l, u):
        l, u = u, l
    m = mediant(l, u)
    return width(l, m), width(m, u), width(l, u)

def fibonacci_pair(k: int):
    a, b = 1, 1
    if k == 0:
        return ERP(1, 1)
    for _ in range(k):
        a, b = a + b, a
    return ERP(a, b)

def run_tests():
    print("lop-ray test")
    for k in range(5):
        x = lop_ray(k)
        y = lop_ray(k + 1)
        print(k, x, y, d_cross(x, y))
        assert x == ERP(k + 1, 1)
        assert d_cross(x, y) == 1

    print("eta-ray test")
    for k in range(6):
        x = eta_ray(k)
        y = fibonacci_pair(k)
        print(k, x, y)
        assert x == y
    for k in range(5):
        older = eta_ray(k)
        newer = eta_ray(k + 1)
        sign = d_cross(older, newer)
        print(k, sign)
        if k % 2 == 0:
            assert sign == 1
        else:
            assert sign == -1

    print("width preservation test")
    w1, w2, w0 = check_width_preservation()
    print(w1, w2, w0)
    assert w1 == w2 == w0

    print("target location and geodesic test")
    a = ERP(2, 3)
    b = ERP(3, 5)
    pa = locate_target(a)
    pb = locate_target(b)
    g = geodesic(a, b)
    print(pa)
    print(pb)
    print(g)
    assert obs_equiv(g[0], a)
    assert obs_equiv(g[-1], b)

    print("structural symbol decomposition test")
    a = ERP(1, 1)
    b = ERP(1, 2)
    c = ERP(1, 3)
    ab = structural_symbol(a, b)
    bc = structural_symbol(b, c)
    ac = structural_symbol(a, c)
    print(ab)
    print(bc)
    print(ac)
    assert len(ac) == len(ab) + len(bc)

    print("barycentric truncation parity test")
    l = ERP(1, 3)
    u = ERP(2, 5)
    print(barycentric_truncation(l, u, 3))
    print(barycentric_truncation(l, u, 4))
    assert barycentric_truncation(l, u, 3) == u
    assert barycentric_truncation(l, u, 4) == l

    print("ZERO-boundary resolution test")
    resolved, empty_trace = zero_resolution(5, 3, 7)
    print(resolved, empty_trace)
    assert resolved == ERP(7, 5)
    assert empty_trace == ERP(0, 3)

    print("koppa test")
    print(koppa(11, 3))
    assert koppa(11, 3) == 2

if __name__ == "__main__":
    run_tests()
\end{lstlisting}

\end{document}